%==============================================================================
\chapter{Beekeeping domain reference and good practices}
%==============================================================================
This chapter specifies the domain knowledge that underpins the management rules,
the configurable thresholds and the system indicators. It constitutes the
functional basis on which the dashboards and the alerts rely.

\section{Honey production factors}
A colony's production depends directly on the weather, pollen availability and
nectar flow. When the nectar flow is high, the queen lays more, the population
peaks and the yield reaches its maximum. In a period of food shortage, laying
and population decrease. The system must therefore correlate the measured
production with periods and seasons, which justifies the filtering of the
dashboards by month, week, season and year.

\section{Location and configuration of a site}
The choice of a site's location determines productivity. The following reference
values can feed the fields and the input checks of the site-management module.
\begin{xltabular}{\textwidth}{>{\bfseries}p{5.5cm} X}
\toprule
\textbf{Criterion} & \textbf{Reference value} \\
\midrule
\endhead
Proximity to resources & Nectar and pollen sources about 1 km away. \\
Access to water & Accessible drinking water, several litres consumed per day and per colony. \\
Distance from dwellings & 100 m in a forest area, 200 m in a shrubland area, 300 m in an open area. \\
Distance from treated crops & At least 3 to 4 km from conventional crops using pesticides. \\
Elevation and inclination & Hives raised about 1.5 m off the ground and slightly tilted to drain moisture. \\
Protection & Shelter from intense sun, wind and floods. \\
\bottomrule
\end{xltabular}

\begin{figure}[htbp]
\centering
\includegraphics[width=0.48\linewidth]{hive-003}
\caption{Example of a hive installed on a site: elevation on legs and landing board.}
\label{fig:ruche-installee}
\end{figure}

\section{Hive types and reference yield}
\begin{xltabular}{\textwidth}{>{\bfseries}p{5.5cm} X}
\toprule
\textbf{Type of hive} & \textbf{Characteristics and yield} \\
\midrule
\endhead
Traditional fixed-comb hive & 6 to 10 kg of comb honey per season, not recommended in organic beekeeping. \\
Top-bar hive & Standard width of 32 to 35 cm, independent inspection of the combs. \\
Frame hive (Langstroth, East-African) & Honey combs reusable for several seasons after harvest. \\
\bottomrule
\end{xltabular}
These orders of magnitude serve as a basis for setting the production thresholds
and estimating the return on investment per hive.

\section{Classification of hive models}
Two large families of hives are distinguished, whose modelling must remain open
to cover local and international uses.
\begin{xltabular}{\textwidth}{>{\bfseries}p{4.5cm} X}
\toprule
\textbf{Family} & \textbf{Representative models} \\
\midrule
\endhead
Frame hives (vertical) & Dadant, Langstroth, Voirnot, Layens, Aumônière, National (GB), WBC, Tonelli. \\
Traditional frameless hives & Warré, Kenyan TBH, Tanzanian TBH, log, straw (skep). \\
Other regional models & Lucitana, Oksman, Smith, CDB, Duvauchelle, Jarry, Congres, Zanders, Slovenian models. \\
\bottomrule
\end{xltabular}
In France, the most common models are the Dadant, Langstroth and Voirnot hives.
All frame hives share a common composition, with dimensions precise to the
millimetre, which reinforces the chosen data model.
\begin{encadre}
\textbf{Model anchor:} the project's reference model is the Dadant hive, whose
body holds up to 10 frames. This value grounds the maximum capacity of 10 frames
per body or per super adopted in the management rules. The system nevertheless
remains configurable to support other models.
\end{encadre}

\subsection{Common composition of a frame hive}
\begin{xltabular}{\textwidth}{>{\bfseries}p{4cm} X}
\toprule
\textbf{Element} & \textbf{Role} \\
\midrule
\endhead
Floor / landing board & Base of the hive and landing zone for the bees. \\
Body (base) & Main brood compartment, holds the base frames. \\
Super & Extension level intended for honey harvesting. \\
Frame & Removable support of the wax combs, installed on an identified slot. \\
Queen excluder & Separates the body from the supers and confines the queen to the body. \\
Crown board and roof & Close and protect the hive from bad weather. \\
Feeder & Allows the food replenishment of the colony. \\
\bottomrule
\end{xltabular}

\begin{figure}[htbp]
\centering
\includegraphics[width=0.85\linewidth]{hive-001}
\caption{Cross-section of a frame hive and role of each element (roof, feeder, crown board, super, queen excluder, body, floor).}
\label{fig:coupe-ruche}
\end{figure}

\section{Colony inspection protocol}
Inspections are preferably carried out on sunny days and focus on precise
checkpoints that structure the visit report.
\begin{itemize}
  \item brood development (eggs, larvae, capped cells)
  \item filling of the cells with honey and pollen
  \item presence of parasites or diseases
  \item collection of nectar, pollen and propolis
\end{itemize}
\begin{encadre}
\textbf{Operational constraint:} a visit must not exceed 45 minutes. Beyond that,
the colony's agitation increases and spreads to the neighbouring hives. The
\emph{duration} field of the visit report can be checked against this reference
limit.
\end{encadre}

\section{Swarming and colony splitting}
An imminent swarming is signalled by the presence of queen cells on the edges of
the combs. A few days before the emergence of a new queen, the old queen leaves
the colony with half the workers. The system must therefore make it possible to
anticipate the split. Three reference methods are provided as operation reasons.
\begin{xltabular}{\textwidth}{>{\bfseries}p{4.5cm} X}
\toprule
\textbf{Method} & \textbf{Principle} \\
\midrule
\endhead
Nucleus formation & Remove 1 to 2 brood combs with bees to a new hive. \\
Queen nucleus & Transfer the queen with a comb free of queen cells. \\
Artificial swarming & Move the mother hive and place a new hive at the old location with 2 to 3 combs of young brood. \\
\bottomrule
\end{xltabular}

\section{Well-being indicators and causes of decline}
A healthy colony presents a well-developed brood nest at the various stages,
cells filled with resources, visible foraging activity and the absence of
parasites or diseases. These criteria feed the qualitative assessment of the
health status and the productivity level rated from 1 to 3 in the visit report.

The absconding or abandonment of a hive results from excessive disturbances or
inadequate conditions (lack of food or water, extreme temperatures). Leaving a
honey reserve to the colony during the harvest reduces this risk. These
phenomena justify the alerts of sudden population or production drop and the
associated configurable thresholds.

\section{Impact on the configurable thresholds}
\begin{xltabular}{\textwidth}{>{\bfseries}p{4.5cm} X}
\toprule
\textbf{System parameter} & \textbf{Business anchor} \\
\midrule
\endhead
Honey-harvest threshold & Reference yield per hive type and per season. \\
Drop-alert threshold & Abnormal drop linked to absconding, disease or food shortage. \\
Maximum visit duration & Limit of 45 minutes per inspection. \\
Site safety distances & 100, 200 or 300 m depending on the area, 3 to 4 km from treated crops. \\
Analysis periods & Correlation of production and season (nectar flow). \\
\bottomrule
\end{xltabular}
